\section{Limitations}\label{sec:limitations}

% source: RM:97, 244; EXP:527-529, 731-732; HER:384-392.
\paragraph{One training run per configuration.}
Each adapter was trained once, with seed 42, because of compute limits. The intervals in this paper therefore cover the variability of sampling evaluation items and groups; they do not cover training variance from initialization and data order, which is substantial in fine-tuning \citep{dodge2020finetuning}. Differences between configurations are conditional on one run each and may not replicate across seeds. The agreement between the pilot and Hermes (within 0.004 overall) is consistency across two bases, not a replication of one configuration.

% source: EXP:737-739; RM:157; TDR:153-154; PRV:91-94; V7:67-69.
\paragraph{The data and the benchmark are generated by the same policy.}
Answers, traces and references are all computed from the policies described in Section~\ref{sec:data}, so the benchmark measures how well a model learned those policies, not whether the policies are good or whether the model reasons in general. Test scenarios come from the same families as the training scenarios, with only the input format held out (deviation D2), so the results speak to transfer across format, not to new scenarios. No completed human review exists for v0.0.6 or v0.0.7; the only calibration evidence for the synthetic labels is v0.0.5's review, which rejected 32 of 127 sampled examples before the policy changes that v0.0.6 made, and 48 examples carry a known labelling quirk.

% source: EVAL:215-217; EXP:561-562, 765-766; RM:84; HER:577-583; cameron2008bootstrap.
\paragraph{Few clusters and uncorrected rows.}
The cluster intervals rest on 61 answerable and 17 should-decline groups, and cluster inference with few clusters tends to over-reject. The 8-item context-prioritization task is not estimable by group. Per-task rows are uncorrected for multiple comparisons, and H1's per-task verdicts were never recomputed by group; the protocol amendment left them under their original item-level rule.

% source: EXP:260-265; RM:99, 247.
\paragraph{No model produced the held-out format.}
\code{format\_valid} is 0.0 for every model and arm. Format is excluded from the score by design, so the results measure judgment on prose answers given JSON prompts; whether a model can also produce the JSON structure the product's parsers would prefer is unmeasured.

% source: EXP:313-315; LMD:500-502; EXP:502-504.
\paragraph{Abstention is a family shortcut, and the decline labels were unlearnable for three of the four runs.}
Section~\ref{sec:analysis} shows that the v0.0.6 adapters decline by scenario family, not by evidence, and whether Logos v0.0.2 does better was not analyzed. For the v0.0.6 runs, 78 of 349 items expect labels absent from training (deviation D3), which caps their scores and is why the cross-model comparisons use the answerable subset. Beyond the decline labels, two families are labelled \code{impact} in the test split and only \code{evidence} in training, so 94 of the 349 test items (26.9\%) need a family--label pairing that training never shows.

% source: EXP:445, 449-451, 674-676; RM:251.
\paragraph{Arms not run.}
The bare base model and the fine-tuned-plus-orchestration arm were never evaluated in full, so H7 is only partly measured, and the comparisons use the orchestrated baseline rather than the registered bare one (deviation D1). Logos v0.0.2 has no base arm, so the effect of fine-tuning on the Reasoning checkpoint is not measured; H9 compares two fine-tuned models only.

% source: EXP:718-739; LMD:487-490; RM:98.
\paragraph{H9 confounds three changes.}
The base checkpoint, the training data and the generation budget changed together, the hardware and torch build differ, and Hermes was not retrained on v0.0.7. The interval's lower bound, $+0.0040$, is close to zero and rests on 61 clusters, where cluster inference tends to over-reject; the verdict is one run, one seed, one benchmark.

% source: EXP:727-730, 710-715, 785-786; L1:321-327; HER:391-392; L1:195-196.
\paragraph{Decoding, selection and the base's budget.}
The Reasoning model was decoded greedily although its card recommends sampling at temperature 1; 6 of 349 completions looped and were scored as generated. Its adapter was selected on a loss that covers the trace, and no answer-only loss exists. For Logos v0.0.1, 41 of the base's answers were cut at 512 tokens, which may inflate that run's base-to-adapter gains. Validation losses are not comparable across tokenizers, so the curves in Figure~\ref{fig:validation} compare in shape only.

% source: RM:250; LMD:509-511; SVR:197-208; EXP:939-968; D12 EXP:920; EXP:733-736; TRN:43-44; RM:254; RM:252.
\paragraph{Serving, provenance and tuning.}
The served GPU differs from the evaluated one; Logos reproduced 8 of 9 checks and its live outputs have not been re-measured on the benchmark. The pilot's base revision is not recoverable, so that run is reproducible only modulo upstream not having moved; the Kaggle image was not pinned (deviation D12) and its torch version is recorded only privately. The hyperparameters are engineering defaults, untuned. The faithfulness measures are text heuristics, not entailment checks. A 24B vision-language configuration in the repository was never trained.

\subsection{Data provenance and licensing}
% source: PRV:42-44; RM:334; LMD:35; EXP:91; RM:230; TPE:20-25.
Every training and test example is synthetic, generated from a catalog of fictional entities and captured through an offline mock backend; no release contains conversations, records or other data from real users, and the pipeline's provenance gate rejects any example from a production-observation lane. The two production bases, Mistral-Nemo-Instruct-2407 and Ministral-3-14B-Reasoning-2512, are released under Apache-2.0; the pilot's base, Ministral-8B-Instruct-2410, is under the Mistral AI Research License, which is why that adapter is research-only and was never served. The trained adapters are distributed privately through the deployment packages rather than published; the training and evaluation code, the frozen records and the data pipeline are public under MIT. The misuse surface of a model that ranks a student's own items and declines on weak evidence is narrow, and the serving endpoints require a key and run under a per-account quota.
